\section{Scalar quotients and the circle branches}
\begin{lemma}[A scalar-orbit test]\label{scalarorbits}
Let $F$ be a finite field and $C\le F^*$. Suppose
$F^+\rtimes C\le H\le A\Gamma L_1(F)$. The affine action of $H$
is two-transitive exactly when its zero stabilizer is transitive on
$F^*/C$. If $\F_p[C]=F$, the action is primitive.
\end{lemma}
\begin{proof}
Translations reduce ordered-pair transitivity to transitivity of the zero
stabilizer on $F^*$. Scalar multiplication by $C$ is transitive on each
of its multiplicative cosets; these cosets are preserved as a system by
all scalar-Frobenius maps. This proves the first equivalence. A block
through zero is an additive subspace invariant under the zero stabilizer.
If it contains a nonzero $v$ and is invariant under $C$, it contains
$\F_p[C]v=F$. Thus no proper nonzero block exists.
\end{proof}

\begin{proposition}[Projective-line circles]\label{projectiveline}
Let $\operatorname{PSL}_2(q)\le G\le\operatorname{P\Gamma L}_2(q)$
in its natural degree-$q+1$ action. Its circle is always residually
weakly primitive. It is locally two-transitive exactly when the induced
action of $G$ on the determinant square classes is transitive. For even
$q$ this condition is automatic; for odd $q$ it says that some element
has nonsquare determinant in its semilinear projective representative.
\end{proposition}
\begin{proof}
The natural special-linear action is two-transitive. Fixing infinity
identifies its remaining points with $\F_q$. Its induced stabilizer
is the affine group $\F_q^+\rtimes C$, where
$C=(\F_q^*)^{\gcd(2,q-1)}$: the diagonal matrix
$\operatorname{diag}(a,a^{-1})$ gives multiplication by $a^2$,
and upper unipotents give every translation. General semilinear
stabilizers give every allowed scalar-Frobenius map. The quotient
character is determinant modulo squares, invariant under projective
scalar changes. Since the special-linear group is transitive, the full
point stabilizer maps onto the outer quotient. Lemma~\ref{scalarorbits}
therefore gives the claimed two-transitivity criterion on the remaining
points, equivalent to three-transitivity on the whole set.

For primitivity, if $q$ is prime, the additive group already has prime
degree. Otherwise every proper subfield has size at most $\sqrt q$,
whereas $|C|\ge(q-1)/2>\sqrt q-1$. Thus $C$ cannot be contained in a
proper subfield and generates $\F_q$; Lemma~\ref{scalarorbits} applies.
Together with Proposition~\ref{circles} this proves the assertion.
The criterion retains mixed diagonal-field extensions; it does not
incorrectly require the acting group to contain $\operatorname{PGL}_2(q)$.
\end{proof}

\begin{lemma}[A characteristic subgroup in a regular normal subgroup]
If a transitive group on $X$ contains a regular normal subgroup $U$
with a nontrivial proper characteristic subgroup $Z$, then it is
imprimitive.
\end{lemma}
\begin{proof}
The $Z$-orbits are blocks, since $Z$ is normal in the full group.
Regularity identifies their size with $|Z|$, strictly between one and
$|X|=|U|$.
\end{proof}

This proves the local exclusion for the natural Suzuki and Ree circle
actions from their precise retained rank-one structure: a point
stabilizer has a regular unipotent radical of order $q^2$ or $q^3$,
respectively, and a proper nontrivial characteristic center. Field and
torus automorphisms preserve this center. Thus the remaining-point
action of the point stabilizer is imprimitive, also for its allowed
outer extensions. No assertion that a Suzuki ovoid is itself a thick
linear space is needed. Natural unitary circles instead preserve the
actual thick unital, so Proposition~\ref{circles} directly excludes them.
Projective actions of dimension at least two preserve their thick
projective lines and are excluded in the same way.

For affine circles over an odd prime field, the actual affine lines
are thick and again give an invariant pencil partition. For extension
field actions on an even field larger than two, the same exclusion
uses the actual invariant extension-field lines, after establishing
their normalizer as in Lemma~\ref{field}. Over $\F_2$, by contrast,
affine lines have only two points and cannot be used for this exclusion.
The natural special-linear family gives three-transitivity, and the
symplectic family gives primitivity but not two-transitivity on the
nonzero vectors by Proposition~\ref{symplectic}. The exceptional
$G_2(2)'$ degree-63 local action requires its own small case: it must
not be excluded by pretending $\operatorname{AG}(6,2)$ is thick.
The following Proposition~\ref{binaryg2} handles it using the natural
module as a retained input and a complete small maximal-subgroup
calculation; the other exceptional action interfaces are kept separate.

The one-dimensional circle branch is covered by Theorem~\ref{spread}
with $\ell=1$: the pencil size $p^d-1$ must be prime. This includes,
for example, $\operatorname{AGL}_1(8)$ with primitive cyclic pencil
of degree seven. It is omitted from the literal stand-alone list in
source Proposition 4.3, whose proof defers the one-dimensional case
to Section 5. The source's final Theorem 6.1 does include it in its
first branch. The corresponding locally two-transitive circle is
only the four-point case among one-dimensional spread actions.

For the exceptional almost-simple branches we will retain exact
permutation-action information, rather than infer primitivity from
abstract group names or orders alone. In particular the $A_7$ action
on fifteen points preserves the lines of $\operatorname{PG}(3,2)$.
It is two-transitive on points but its point stabilizer preserves
seven pairs on the other fourteen points. Its circle is not locally
primitive and cannot be three-transitive. This is the previously
accepted Topic G circle correction, not a new conjecture result.

\begin{proposition}[The binary exceptional local action]\label{binaryg2}
Retain the natural transitive binary representation of dimension six of
$N=G_2(2)'\cong\operatorname{PSU}_3(3)$. Every group
$N\le H\le N_{\GL_6(2)}(N)$ is primitive but not two-transitive on
the sixty-three nonzero vectors. Thus $2^6\rtimes H$ has a residually
weakly primitive circle action and is not locally two-transitive.
\end{proposition}
\begin{proof}
The small exact finite premise supplied by \texttt{small\_g2\_premise.g}
is the complete maximal-subgroup order list
\[
 |N|=6048,\qquad 96,96,168,216,\qquad |\Aut(N)|=12096.
\]
It is computed in GAP from $\operatorname{PSU}(3,3)$, including a
simplicity check and both actual degree-63 coset actions. This is a
required finite premise, not an assertion derived solely from a subgroup
order. In the natural transitive action a vector stabilizer has order
$6048/63=96$. Any proper overgroup of it lies in a maximal subgroup,
and no order in the displayed complete list is a larger multiple of
96. Thus the stabilizer is maximal and the action is primitive.

For the normalizer bound we can avoid assuming an unverified full
matrix normalizer. Transitivity on nonzero vectors implies irreducibility.
Its commuting algebra is a finite division ring, hence a field $\F_{2^e}$,
with $e\mid6$. If $e=3$ or $6$, Lagrange's theorem would embed $N$
in $\GL_2(8)$ or $\GL_1(64)$, whose orders are not divisible by
6048. If $e=2$, simplicity and perfectness put $N$ inside
$\operatorname{SL}_3(4)$; its intersection with the central subgroup of order three
is trivial, so it embeds in $\operatorname{PSL}_3(4)$, of order 20160, again not
divisible by 6048. Consequently $e=1$, and its invertible commuting
algebra is trivial. Conjugation embeds the entire linear normalizer
in $\Aut(N)$. A two-transitive action on sixty-three points would
require order divisible by $63\cdot62$, which does not divide 12096.
Every intermediate group inherits primitivity from $N$, but none
can be two-transitive. Proposition~\ref{circles} proves the conclusion.
\end{proof}

\begin{proposition}[Orthogonal-point symplectic actions]\label{orthogonalcircles}
Retain the standard two-transitive actions of
$\operatorname{Sp}_{2n}(2)$ of degrees
$2^{2n-1}\pm2^{n-1}$, $n\geq3$, and the classical rank-three subdegrees
of their orthogonal point stabilizers. Both circle actions are locally
primitive and neither is locally two-transitive.
\end{proposition}
\begin{proof}
Put $b=2^{2n-2}$ and $d=2^{2n-1}\pm2^{n-1}-1$.
On the $d$ remaining points, a point stabilizer of the orthogonal group
has orbits of lengths $1$, $a=d-1-b$, and $b$; these are the precise
retained rank-three data. A block containing the fixed point is a union
of these orbits, so its only possible proper nontrivial sizes are
$a+1=d-b$ and $b+1$. The first is an odd integer greater than one; if it
divided $d$, it would divide the power of two $b$, impossible. For the
minus sign, $d/2<b+1<d$. For the plus sign,
$2(b+1)<d<3(b+1)$, since $n\geq3$. Thus $b+1$ cannot divide $d$ either.
The local action is primitive. Its rank is three, so it is not
two-transitive. The circle criterion proves both conclusions.
\end{proof}
